% ═══════════════════════════════════════════════════════════════
%  chapters/ch17-appimage.tex
%  AppImage: упаковка приложения для Linux
% ═══════════════════════════════════════════════════════════════

\chapter{AppImage: упаковка приложения}
\label{ch:appimage}

\epigraph{<<AppImage~--- это способ сказать пользователю:
<<Запусти~--- и всё заработает>>. Без установки, без зависимостей>>.}{}

\section{Введение}
\label{sec:appimage-intro}

В \texttt{ch16} мы \emph{отладили} приложение. Оно \emph{работает}
\emph{на} \emph{нашей} \emph{машине}. Но \emph{что}, \emph{если}
\emph{пользователь} \emph{хочет} \emph{запустить} \emph{его} \emph{на}
\emph{своей} \emph{системе}? \emph{У} \emph{него} \emph{может} \emph{не}
\emph{быть} \emph{GTK4}, \emph{GLib} \emph{нужной} \emph{версии},
\emph{правильных} \emph{библиотек}.

\textbf{Решение}~--- \textbf{AppImage}. \emph{Один} \emph{файл},
\emph{который} \emph{содержит} \emph{всё}: \emph{бинарник},
\emph{библиотеки}, \emph{иконки}. \emph{Запустил}~---
\emph{работает}.

\begin{ideabox}
\textbf{AppImage~--- это \emph{контейнер}.} Как \emph{чемодан}:
\emph{внутри}~--- \emph{всё} \emph{необходимое}. \emph{Пользователь}
\emph{не} \emph{думает} \emph{о} \emph{зависимостях}. \emph{Он}
\emph{просто} \emph{запускает}.
\end{ideabox}

\subsection{Что вы узнаете}

\begin{itemize}
    \item Что такое \textbf{AppImage} и \emph{зачем} он нужен.
    \item Как \emph{собрать} AppImage через \texttt{linuxdeploy}.
    \item Как \emph{добавить} иконку и \texttt{.desktop}-файл.
    \item Как \emph{проверить} AppImage на \emph{чистой} системе.
\end{itemize}

\section{Что такое AppImage}
\label{sec:appimage-what}

\textbf{AppImage}~--- это \emph{формат} упаковки приложений
для Linux. \emph{Один} \emph{файл}, \emph{который}:
\begin{enumerate}
    \item \textbf{Содержит} \emph{все} \emph{библиотеки}.
    \item \textbf{Монтируется} \emph{при} \emph{запуске}.
    \item \textbf{Работает} \emph{на} \emph{любом} \emph{дистрибутиве}.
    \item \textbf{Не} \emph{требует} \emph{установки}.
\end{enumerate}

\begin{table}[htbp]
\centering
\caption{Сравнение форматов}
\label{tab:appimage-vs}
\small
\begin{tabularx}{\textwidth}{@{}l L L@{}}
\toprule
\textbf{Формат} & \textbf{Плюсы} & \textbf{Минусы} \\
\midrule
\textbf{AppImage} & Один файл, без установки & \emph{Большой} \emph{размер} \\
\textbf{deb/rpm}  & Малый размер & Зависит от \emph{дистрибутива} \\
\textbf{Flatpak}  & Песочница & Сложная \emph{настройка} \\
\textbf{Snap}     & Автообновление & Нужен \emph{демон} \\
\bottomrule
\end{tabularx}
\end{table}

\begin{keybox}
\textbf{AppImage~--- для \emph{пользователя}.} Он \emph{скачал},
\emph{сделал} \emph{chmod +x}, \emph{запустил}. \emph{Всё}.
\emph{Никаких} \emph{команд}, \emph{никаких} \emph{зависимостей}.
\end{keybox}

\section{Сборка AppImage}
\label{sec:appimage-build}

\textbf{Инструмент}~--- \texttt{linuxdeploy}. \emph{Он}
\emph{собирает} \emph{AppImage} \emph{из} \emph{собранного}
\emph{бинарника}.

\subsection{Установка}

\begin{listing}
\begin{minted}[frame=none, bgcolor=gray!5]{bash}
# Скачиваем linuxdeploy
wget https://github.com/linuxdeploy/linuxdeploy/releases/\
     download/continuous/linuxdeploy-x86_64.AppImage
chmod +x linuxdeploy-x86_64.AppImage
sudo mv linuxdeploy-x86_64.AppImage /usr/local/bin/linuxdeploy

# Плагин для GTK
wget https://github.com/linuxdeploy/linuxdeploy-plugin-gtk/\
     releases/download/continuous/\
     linuxdeploy-plugin-gtk.sh
chmod +x linuxdeploy-plugin-gtk.sh
sudo mv linuxdeploy-plugin-gtk.sh /usr/local/bin/
\end{minted}
\caption{Установка \texttt{linuxdeploy}}
\label{lst:appimage-install}
\end{listing}

\subsection{Сборка}

\textbf{Три} \emph{шага}:

\begin{enumerate}
    \item \textbf{Собрать} \emph{бинарник}.
    \item \textbf{Подготовить} \emph{AppDir} (\emph{директорию}
          \emph{с} \emph{приложением}).
    \item \textbf{Собрать} \emph{AppImage}.
\end{enumerate}

\begin{listing}[htbp]
\begin{minted}[frame=none, bgcolor=gray!5]{bash}
# 1. Собираем бинарник
cd pc-app
make clean
make

# 2. Готовим AppDir
mkdir -p AppDir/usr/bin
mkdir -p AppDir/usr/share/applications
mkdir -p AppDir/usr/share/icons/hicolor/256x256/apps

cp osc_gen_ui AppDir/usr/bin/

# 3. Desktop-файл
cat > AppDir/usr/share/applications/osc_gen_ui.desktop << 'EOF'
[Desktop Entry]
Name=OscGen UI
Exec=osc_gen_ui
Icon=osc_gen_ui
Type=Application
Categories=Development;Utility;
EOF

# 4. Иконка (если есть)
# cp icon.png AppDir/usr/share/icons/hicolor/256x256/apps/osc_gen_ui.png

# 5. Собираем AppImage
linuxdeploy --appdir AppDir \
    --executable AppDir/usr/bin/osc_gen_ui \
    --desktop-file AppDir/usr/share/applications/osc_gen_ui.desktop \
    --icon-file AppDir/usr/share/icons/hicolor/256x256/apps/osc_gen_ui.png

appimagetool AppDir
\end{minted}
\caption{Сборка AppImage}
\label{lst:appimage-build}
\end{listing}

\begin{keybox}
\textbf{\texttt{linuxdeploy}~--- умный.} Он \emph{сам} \emph{находит}
\emph{зависимости} (\emph{GTK}, \emph{GLib}, \emph{Cairo}) и
\emph{копирует} их \emph{в} \emph{AppDir}.
\emph{Пользователь} \emph{не} \emph{думает} \emph{о} \emph{них}.
\end{keybox}

\begin{warnbox}
\textbf{GTK-плагин обязателен.} Без него \emph{тема},
\emph{шрифты} и \emph{иконки} \emph{не} \emph{работают}.
\emph{Плагин} \emph{копирует} \emph{все} \emph{темы} \emph{и}
\emph{шрифты} \emph{в} \emph{AppDir}.
\end{warnbox}

\section{Проверка}
\label{sec:appimage-test}

\textbf{Как} \emph{убедиться}, \emph{что} \emph{AppImage}
\emph{работает}?

\begin{enumerate}
    \item \textbf{Скопировать} \emph{на} \emph{другую} \emph{машину}
          (\emph{или} \emph{в} \emph{Docker}).
    \item \textbf{Запустить} \emph{без} \emph{установки}.
    \item \textbf{Проверить}, \emph{что} \emph{всё} \emph{работает}.
\end{enumerate}

\begin{minted}[frame=none, bgcolor=gray!5]{bash}
# На чистой системе
chmod +x OscGen_UI-x86_64.AppImage
./OscGen_UI-x86_64.AppImage
\end{minted}

\begin{ideabox}
\textbf{Если} \emph{AppImage} \emph{запускается} \emph{на}
\emph{чистой} \emph{системе}~--- \emph{всё} \emph{хорошо}.
\emph{Если} \emph{нет}~--- \emph{значит}, \emph{linuxdeploy}
\emph{пропустил} \emph{какую-то} \emph{библиотеку}.
\emph{Смотрим} \emph{вывод} \emph{команды}.
\end{ideabox}

\section{Уроки главы}
\label{sec:appimage-lessons}

\begin{enumerate}
    \item \textbf{AppImage~--- один файл.} Без установки.
    \item \textbf{linuxdeploy} \emph{собирает} \emph{AppImage}.
    \item \textbf{GTK-плагин} \emph{копирует} \emph{темы}
          \emph{и} \emph{шрифты}.
    \item \textbf{Проверка} \emph{на} \emph{чистой} \emph{системе}
          \emph{обязательна}.
    \item \textbf{Размер} \emph{AppImage}~--- \emph{большой}
          (\emph{100}--\emph{300} \emph{МБ}).
\end{enumerate}

\section{Чек-лист грабель}
\label{sec:appimage-pitfalls}

\begin{table}[htbp]
\centering
\caption{Грабли AppImage}
\label{tab:appimage-pitfalls}
\small
\begin{tabularx}{\textwidth}{@{}c L L@{}}
\toprule
\textbf{№} & \textbf{Грабля} & \textbf{Решение} \\
\midrule
1 & AppImage не запускается & \emph{Проверить} \emph{зависимости} \\
2 & Нет темы GTK & \emph{GTK-плагин} \\
3 & Нет иконки & \emph{Добавить} \emph{в} \emph{AppDir} \\
4 & Нет \texttt{.desktop} & \emph{Создать} \emph{файл} \\
5 & Большой размер & \emph{Оптимизировать} \emph{библиотеки} \\
6 & \emph{Нет} \emph{прав} \emph{на} \emph{запуск} & \texttt{chmod +x} \\
7 & Ошибка \emph{сегментации} & \emph{Проверить} \emph{на} \emph{чистой} \emph{системе} \\
8 & \emph{Не} \emph{видит} \emph{порт} & \emph{Правила} \emph{udev} \\
\bottomrule
\end{tabularx}
\end{table}

\section{Упражнения}
\label{sec:appimage-exercises}

\begin{exercisebox}
\textbf{Упражнение 1.} Соберите \emph{AppImage} для
\emph{вашего} \emph{приложения}. \emph{Какой}
\emph{размер}?
\end{exercisebox}

\begin{exercisebox}
\textbf{Упражнение 2.} Добавьте \emph{иконку}.
\emph{Как} \emph{проверить}, \emph{что} \emph{она}
\emph{отображается}?
\end{exercisebox}

\needspace{5\baselineskip}
\begin{exercisebox}
\textbf{Упражнение 3.} Запустите \emph{AppImage} в
\emph{Docker}-контейнере (\emph{чистый} \emph{Ubuntu}).
\emph{Работает}?
\end{exercisebox}

\begin{exercisebox}
\textbf{Упражнение 4.} Предложите \emph{способ}
\emph{уменьшить} \emph{размер} \emph{AppImage}.
\emph{Что} \emph{можно} \emph{исключить}?
\end{exercisebox}

\section{Проверьте сами}
\label{sec:appimage-check}

\begin{enumerate}
    \item Соберите \emph{бинарник}.
    \item Подготовьте \emph{AppDir}.
    \item Запустите \emph{linuxdeploy}.
    \item Соберите \emph{AppImage} \emph{через}
          \emph{appimagetool}.
    \item Скопируйте \emph{на} \emph{другую} \emph{машину}.
    \item \emph{Запустите}. \emph{Проверьте}, \emph{что}
          \emph{всё} \emph{работает}.
    \item \emph{Сделайте} \emph{скриншот}~--- \emph{пришлите}
          \emph{в} \emph{книгу}.
\end{enumerate}

\begin{ideabox}
\textbf{Если AppImage запускается~--- вы закончили.}
Это \emph{главное}. \emph{Дальше}~--- \emph{приложения}:
\emph{исходный} \emph{код}, \emph{сборка}, \emph{глоссарий}.
\emph{Книга} \emph{почти} \emph{готова}.
\end{ideabox}


